% Part: lambda-calculus
% Chapter: introduction
% Section: basic-pr-lambda

\documentclass[../../../include/open-logic-section]{subfiles}

\begin{document}

\olfileid{lam}{int}{bas}
\olsection{بنسټيزې بازګشتي تابعې \usetoken{S}{lambda definable} دي}

\begin{lem}
تابعې $\Zero$، $\Succ$ او $\Proj{n}{i}$ !!{lambda definable} دي۔
\end{lem}

\begin{proof}
$\Zero$ يوازې $\lambd[u][\lambd[x][\lambd[y][y]]]$ دے۔
\textbf{د ژباړې د نسخې سمون (OLLAM-004):} سرچينې د صفر د يوځاييزې
تابع دپاره يوازې د چرچ د صفر عددنښه راوړې وه۔ دلته د تابعې د ننوت
دپاره يو بهرنی تړون زيات شوے دے، څو هر ننوت هماغه د صفر عددنښه
ورکړي؛ د همدې کتاب د لامبډا-تعريفېدنې وروستے څپرکے هم دا بهرنی
تړون په ښکاره کاروي۔

د تالي تابع~$\Succ$ تعريف دا دے:
$\fn{Succ}(u) = \lambd[x][\lambd[y][x(uxy)]]$۔ ځير شئ چې دا
ولې کار کوي: هره عددنښه~$\num{n}$ د تکراروونکي په توګه وګڼئ۔ د
هرې تابعې~$f$ دپاره $\fn{Succ}(\num{n},f)$ داسې تابع ده چې پر
ننوت~$y$، لومړے له~$y$ څخه په پيل سره~$f$،~$n$ ځله تطبيقوي او
بيا ئې يو ځل نور هم تطبيقوي۔

د پروجيکشنونو په هکله نور څۀ نشته: $\fn{Proj}^n_i(x_0, \dots,
x_{n-1}) = x_i$۔ په بله وينا، زموږ د نښيزو دودونو له مخې
$\fn{Proj}^n_i$ لامبډا ترم
$\lambd[x_0][\dots \lambd[x_{n-1}][x_i]]$ دے۔
\end{proof}

\end{document}
